%% ==========================================================================
%%  SECTION 1 --- INTRODUCTION
%%
%%  Job: make the reader want the rest of the survey. Establish the problem,
%%  the two halves, and the shape of the argument. Roughly 1.5 pages.
%% ==========================================================================
\section{Introduction}
\label{sec:intro}

\todo{Open with the concrete harm, not with a definition. \pirsonaplain{} \S1 does
      this well --- consumption data is invaluable to ``unscrupulous marketers,
      identity thieves\dots{} and hostile foreign agents''. Borrow the framing,
      not the words.}

\subsection{Two halves of one problem}
\label{sec:intro:two-halves}

\todo{The central framing device of the whole survey. A recommender that
      protects its users must do two separate things:
      (i) TRAIN a model on private ratings without learning them, and
      (ii) DELIVER the recommended item without learning which item.
      Doing only (i) is not enough --- the fetch betrays the recommendation.
      Doing only (ii) is not enough --- the model was built from exposed data.
      Establish this here and refer back to it in every later section.}

\subsection{The apparent contradiction}
\label{sec:intro:contradiction}

\todo{Why this looked impossible: PIR hides WHICH item a user fetches, while
      collaborative filtering works precisely by relating users to each other
      through what they fetched. \pirsonaplain{} \S1 states this tension
      explicitly --- quote it, then say it has been resolved twice over, in
      two different ways, by two literatures that did not talk to each other.}

\subsection{Scope of this survey}
\label{sec:intro:scope}

\todo{State what is in and what is out.
      IN:  MPC- and PIR-based approaches; matrix factorisation; private
           nearest-neighbour retrieval; the leakage that survives.
      OUT: pure differential privacy without a cryptographic layer; trusted
           hardware (mentioned in \S2 for contrast only); neural recommenders.
      Say why the cuts were made --- an unjustified scope cut reads as an
      unread paper.}

\subsection{Contributions of this survey}
\label{sec:intro:contributions}

\todo{Be honest: a survey's contribution is a map and a gap, not a result.
      Claim exactly three things:
      (1) the first side-by-side treatment of the private-training and
          private-retrieval literatures for recommendation;
      (2) identification of a concrete, author-acknowledged gap between them;
      (3) two cheap, checkable research questions arising from it (N1, N2 ---
          see \Cref{sec:gap}).}

\subsection{Roadmap}
\label{sec:intro:roadmap}

\todo{One short paragraph mapping \S2--\S9. Make the arc audible: primitives,
      then training, then retrieval, then who tried both, then the gap.}
